
\documentclass{article}
\usepackage{colortbl}
\usepackage{makecell}
\usepackage{multirow}
\usepackage{supertabular}

\begin{document}

\newcounter{utterance}

\centering \large Interaction Transcript for game `matchit\_ascii', experiment `different\_grid', episode 6 with Qwen3.6{-}35B{-}A3B{-}FP8{-}without{-}reasoning.
\vspace{24pt}

{ \footnotesize  \setcounter{utterance}{1}
\setlength{\tabcolsep}{0pt}
\begin{supertabular}{c@{$\;$}|p{.15\linewidth}@{}p{.15\linewidth}p{.15\linewidth}p{.15\linewidth}p{.15\linewidth}p{.15\linewidth}}
    \# & $\;$A & \multicolumn{4}{c}{Game Master} & $\;\:$B\\
    \hline

    \theutterance \stepcounter{utterance}  
    & & \multicolumn{4}{p{0.6\linewidth}}{
        \cellcolor[rgb]{0.9,0.9,0.9}{
            \makecell[{{p{\linewidth}}}]{
                \texttt{\tiny{[P1$\langle$GM]}}
                \texttt{Te osalete ühises arvamismängus.} \\
\texttt{Eesmärk on välja selgitada, kas see ruudustik ja teine pilt, mida ainult mina näen, on samad.} \\
\texttt{Palun kirjeldage kõigepealt oma ruudustikku. Seejärel annan mina oma kirjelduse ja me saame teineteisele esitada küsimusi ruudustike kohta, et välja selgitada, kas need on samad. Alustage nüüd oma lühikest ruudustiku kirjeldust märgiga "KIRJELDUS:", millele järgneb kirjeldus. Ärge lisage midagi muud. See on teie ruudustik:} \\
\texttt{X ▢ ▢ ▢ X} \\
\texttt{X ▢ ▢ ▢ X} \\
\texttt{X X X X X} \\
\texttt{X ▢ ▢ ▢ X} \\
\texttt{X ▢ ▢ ▢ X} \\
            }
        }
    }
    & \\ \\

    \theutterance \stepcounter{utterance}  
    & \multicolumn{4}{p{0.6\linewidth}}{
        \cellcolor[rgb]{0.9,0.9,0.9}{
            \makecell[{{p{\linewidth}}}]{
                \texttt{\tiny{[P1$\rangle$GM]}}
                \texttt{KIRJELDUS: Ruudustik on 5x5. Kolmas rida koosneb täielikult X{-}dest. Esimesed ja viimased veerud (1. ja 5. veerg) sisaldavad kõigis ridades X{-}e. Keskmine 3x3 ala (ridu 1{-}5, veerge 2{-}4) koosneb tühjadest ruutudest (▢).} \\
            }
        }
    }
    & & \\ \\

    \theutterance \stepcounter{utterance}  
    & & & \multicolumn{2}{p{0.3\linewidth}}{
        \cellcolor[rgb]{0.9,0.9,0.9}{
            \makecell[{{p{\linewidth}}}]{
                \texttt{\tiny{[GM$|$GM]}}
                \texttt{continue} \\
            }
        }
    }
    & & \\ \\

    \theutterance \stepcounter{utterance}  
    & & \multicolumn{4}{p{0.6\linewidth}}{
        \cellcolor[rgb]{0.9,0.9,0.9}{
            \makecell[{{p{\linewidth}}}]{
                \texttt{\tiny{[GM$\rangle$P2]}}
                \texttt{Te osalete ühises arvamismängus.} \\
\texttt{Eesmärk on välja selgitada, kas see ruudustik ja teine pilt, mida ainult mina näen, on samad.} \\
\texttt{Palun kirjeldage kõigepealt oma ruudustikku. Seejärel annan mina oma kirjelduse ja me saame teineteisele esitada küsimusi ruudustike kohta, et välja selgitada, kas need on samad. Alustage nüüd oma lühikest ruudustiku kirjeldust märgiga "KIRJELDUS:", millele järgneb kirjeldus. Ärge lisage midagi muud. See on teie ruudustik:} \\
\texttt{▢ ▢ ▢ ▢ ▢} \\
\texttt{X X X ▢ X} \\
\texttt{X X ▢ X X} \\
\texttt{X ▢ X X X} \\
\texttt{▢ ▢ ▢ ▢ ▢} \\
            }
        }
    }
    & \\ \\

    \theutterance \stepcounter{utterance}  
    & & & \multicolumn{4}{p{0.6\linewidth}}{
        \cellcolor[rgb]{0.9,0.9,0.9}{
            \makecell[{{p{\linewidth}}}]{
                \texttt{\tiny{[GM$\langle$P2]}}
                \texttt{KIRJELDUS: 5x5 ruudustik, kus ääred on tühjad (▢) ja keskosas on X{-}tähtedega täidetud ruudud, moodustades umbes ristkülikukujulise mustri, kus nurkades on tühjad ruudud.} \\
            }
        }
    }
     \\ \\

    \theutterance \stepcounter{utterance}  
    & & & \multicolumn{2}{p{0.3\linewidth}}{
        \cellcolor[rgb]{0.9,0.9,0.9}{
            \makecell[{{p{\linewidth}}}]{
                \texttt{\tiny{[GM$|$GM]}}
                \texttt{continue} \\
            }
        }
    }
    & & \\ \\

    \theutterance \stepcounter{utterance}  
    & & \multicolumn{4}{p{0.6\linewidth}}{
        \cellcolor[rgb]{0.9,0.9,0.9}{
            \makecell[{{p{\linewidth}}}]{
                \texttt{\tiny{[GM$\rangle$P2]}}
                \texttt{See on minu KIRJELDUSKIRJELDUS: Ruudustik on 5x5. Kolmas rida koosneb täielikult X{-}dest. Esimesed ja viimased veerud (1. ja 5. veerg) sisaldavad kõigis ridades X{-}e. Keskmine 3x3 ala (ridu 1{-}5, veerge 2{-}4) koosneb tühjadest ruutudest (▢).} \\
\texttt{Nüüd esitage küsimus, et teada saada minu ruudustiku uusi aspekte, mis võivad erineda teie ruudustikust. Alustage "KÜSIMUS:" ja ärge lisage midagi muud.} \\
            }
        }
    }
    & \\ \\

    \theutterance \stepcounter{utterance}  
    & & & \multicolumn{4}{p{0.6\linewidth}}{
        \cellcolor[rgb]{0.9,0.9,0.9}{
            \makecell[{{p{\linewidth}}}]{
                \texttt{\tiny{[GM$\langle$P2]}}
                \texttt{KÜSIMUS: Kas teie ruudustiku esimene rida sisaldab X{-}e või on see täielikult tühi?} \\
            }
        }
    }
     \\ \\

    \theutterance \stepcounter{utterance}  
    & & & \multicolumn{2}{p{0.3\linewidth}}{
        \cellcolor[rgb]{0.9,0.9,0.9}{
            \makecell[{{p{\linewidth}}}]{
                \texttt{\tiny{[GM$|$GM]}}
                \texttt{continue} \\
            }
        }
    }
    & & \\ \\

    \theutterance \stepcounter{utterance}  
    & & \multicolumn{4}{p{0.6\linewidth}}{
        \cellcolor[rgb]{0.9,0.9,0.9}{
            \makecell[{{p{\linewidth}}}]{
                \texttt{\tiny{[P1$\langle$GM]}}
                \texttt{See on minu KIRJELDUSKIRJELDUS: 5x5 ruudustik, kus ääred on tühjad (▢) ja keskosas on X{-}tähtedega täidetud ruudud, moodustades umbes ristkülikukujulise mustri, kus nurkades on tühjad ruudud.} \\
\texttt{KÜSIMUS: Kas teie ruudustiku esimene rida sisaldab X{-}e või on see täielikult tühi?Alustage oma vastust "VASTUS:" ja ärge lisage midagi muud.} \\
            }
        }
    }
    & \\ \\

    \theutterance \stepcounter{utterance}  
    & \multicolumn{4}{p{0.6\linewidth}}{
        \cellcolor[rgb]{0.9,0.9,0.9}{
            \makecell[{{p{\linewidth}}}]{
                \texttt{\tiny{[P1$\rangle$GM]}}
                \texttt{VASTUS: Jah, minu ruudustiku esimene rida sisaldab X{-}e. Täpsemalt on esimene ja viimane ruut X{-}ega täidetud, keskmised kolm ruutu on tühjad (▢).} \\
            }
        }
    }
    & & \\ \\

    \theutterance \stepcounter{utterance}  
    & & & \multicolumn{2}{p{0.3\linewidth}}{
        \cellcolor[rgb]{0.9,0.9,0.9}{
            \makecell[{{p{\linewidth}}}]{
                \texttt{\tiny{[GM$|$GM]}}
                \texttt{continue} \\
            }
        }
    }
    & & \\ \\

    \theutterance \stepcounter{utterance}  
    & & \multicolumn{4}{p{0.6\linewidth}}{
        \cellcolor[rgb]{0.9,0.9,0.9}{
            \makecell[{{p{\linewidth}}}]{
                \texttt{\tiny{[P1$\langle$GM]}}
                \texttt{Nüüd esitage küsimus, et teada saada minu ruudustiku uusi aspekte, mis võivad erineda teie ruudustikust. Alustage "KÜSIMUS:" ja ärge lisage midagi muud.} \\
            }
        }
    }
    & \\ \\

    \theutterance \stepcounter{utterance}  
    & \multicolumn{4}{p{0.6\linewidth}}{
        \cellcolor[rgb]{0.9,0.9,0.9}{
            \makecell[{{p{\linewidth}}}]{
                \texttt{\tiny{[P1$\rangle$GM]}}
                \texttt{KÜSIMUS: Kas teie ruudustiku keskosas (ridu 2{-}4, veerge 2{-}4) on X{-}tähed või on need tühjad?} \\
            }
        }
    }
    & & \\ \\

    \theutterance \stepcounter{utterance}  
    & & & \multicolumn{2}{p{0.3\linewidth}}{
        \cellcolor[rgb]{0.9,0.9,0.9}{
            \makecell[{{p{\linewidth}}}]{
                \texttt{\tiny{[GM$|$GM]}}
                \texttt{continue} \\
            }
        }
    }
    & & \\ \\

    \theutterance \stepcounter{utterance}  
    & & \multicolumn{4}{p{0.6\linewidth}}{
        \cellcolor[rgb]{0.9,0.9,0.9}{
            \makecell[{{p{\linewidth}}}]{
                \texttt{\tiny{[GM$\rangle$P2]}}
                \texttt{VASTUS: Jah, minu ruudustiku esimene rida sisaldab X{-}e. Täpsemalt on esimene ja viimane ruut X{-}ega täidetud, keskmised kolm ruutu on tühjad (▢).} \\
\texttt{KÜSIMUS: Kas teie ruudustiku keskosas (ridu 2{-}4, veerge 2{-}4) on X{-}tähed või on need tühjad?Alustage oma vastust "VASTUS:" ja ärge lisage midagi muud.} \\
            }
        }
    }
    & \\ \\

    \theutterance \stepcounter{utterance}  
    & & & \multicolumn{4}{p{0.6\linewidth}}{
        \cellcolor[rgb]{0.9,0.9,0.9}{
            \makecell[{{p{\linewidth}}}]{
                \texttt{\tiny{[GM$\langle$P2]}}
                \texttt{VASTUS: Jah, minu ruudustiku keskosas (ridu 2{-}4, veerge 2{-}4) on X{-}tähed. Täpsemalt: rida 2 on X X ▢ X X, rida 3 on X X ▢ X X, rida 4 on X ▢ X X X.} \\
            }
        }
    }
     \\ \\

    \theutterance \stepcounter{utterance}  
    & & & \multicolumn{2}{p{0.3\linewidth}}{
        \cellcolor[rgb]{0.9,0.9,0.9}{
            \makecell[{{p{\linewidth}}}]{
                \texttt{\tiny{[GM$|$GM]}}
                \texttt{continue} \\
            }
        }
    }
    & & \\ \\

    \theutterance \stepcounter{utterance}  
    & & \multicolumn{4}{p{0.6\linewidth}}{
        \cellcolor[rgb]{0.9,0.9,0.9}{
            \makecell[{{p{\linewidth}}}]{
                \texttt{\tiny{[GM$\rangle$P2]}}
                \texttt{Nüüd esitage küsimus, et teada saada minu ruudustiku uusi aspekte, mis võivad erineda teie ruudustikust. Alustage "KÜSIMUS:" ja ärge lisage midagi muud.} \\
            }
        }
    }
    & \\ \\

    \theutterance \stepcounter{utterance}  
    & & & \multicolumn{4}{p{0.6\linewidth}}{
        \cellcolor[rgb]{0.9,0.9,0.9}{
            \makecell[{{p{\linewidth}}}]{
                \texttt{\tiny{[GM$\langle$P2]}}
                \texttt{KÜSIMUS: Kas teie ruudustiku neljas rida sisaldab X{-}e või on see täielikult tühi?} \\
            }
        }
    }
     \\ \\

    \theutterance \stepcounter{utterance}  
    & & & \multicolumn{2}{p{0.3\linewidth}}{
        \cellcolor[rgb]{0.9,0.9,0.9}{
            \makecell[{{p{\linewidth}}}]{
                \texttt{\tiny{[GM$|$GM]}}
                \texttt{continue} \\
            }
        }
    }
    & & \\ \\

    \theutterance \stepcounter{utterance}  
    & & \multicolumn{4}{p{0.6\linewidth}}{
        \cellcolor[rgb]{0.9,0.9,0.9}{
            \makecell[{{p{\linewidth}}}]{
                \texttt{\tiny{[P1$\langle$GM]}}
                \texttt{VASTUS: Jah, minu ruudustiku keskosas (ridu 2{-}4, veerge 2{-}4) on X{-}tähed. Täpsemalt: rida 2 on X X ▢ X X, rida 3 on X X ▢ X X, rida 4 on X ▢ X X X.} \\
\texttt{KÜSIMUS: Kas teie ruudustiku neljas rida sisaldab X{-}e või on see täielikult tühi?Alustage oma vastust "VASTUS:" ja ärge lisage midagi muud.} \\
            }
        }
    }
    & \\ \\

    \theutterance \stepcounter{utterance}  
    & \multicolumn{4}{p{0.6\linewidth}}{
        \cellcolor[rgb]{0.9,0.9,0.9}{
            \makecell[{{p{\linewidth}}}]{
                \texttt{\tiny{[P1$\rangle$GM]}}
                \texttt{VASTUS: Jah, minu ruudustiku neljas rida sisaldab X{-}e. Esimene ja viimane ruut on X{-}ega täidetud, keskmised kolm ruutu on tühjad (▢).} \\
            }
        }
    }
    & & \\ \\

    \theutterance \stepcounter{utterance}  
    & & & \multicolumn{2}{p{0.3\linewidth}}{
        \cellcolor[rgb]{0.9,0.9,0.9}{
            \makecell[{{p{\linewidth}}}]{
                \texttt{\tiny{[GM$|$GM]}}
                \texttt{continue} \\
            }
        }
    }
    & & \\ \\

    \theutterance \stepcounter{utterance}  
    & & \multicolumn{4}{p{0.6\linewidth}}{
        \cellcolor[rgb]{0.9,0.9,0.9}{
            \makecell[{{p{\linewidth}}}]{
                \texttt{\tiny{[P1$\langle$GM]}}
                \texttt{Nüüd esitage küsimus, et teada saada minu ruudustiku uusi aspekte, mis võivad erineda teie ruudustikust. Alustage "KÜSIMUS:" ja ärge lisage midagi muud.} \\
            }
        }
    }
    & \\ \\

    \theutterance \stepcounter{utterance}  
    & \multicolumn{4}{p{0.6\linewidth}}{
        \cellcolor[rgb]{0.9,0.9,0.9}{
            \makecell[{{p{\linewidth}}}]{
                \texttt{\tiny{[P1$\rangle$GM]}}
                \texttt{KÜSIMUS: Kas teie ruudustiku viies rida on identne esimese reaga (X ▢ ▢ ▢ X) või on seal erinevusi?} \\
            }
        }
    }
    & & \\ \\

    \theutterance \stepcounter{utterance}  
    & & & \multicolumn{2}{p{0.3\linewidth}}{
        \cellcolor[rgb]{0.9,0.9,0.9}{
            \makecell[{{p{\linewidth}}}]{
                \texttt{\tiny{[GM$|$GM]}}
                \texttt{continue} \\
            }
        }
    }
    & & \\ \\

    \theutterance \stepcounter{utterance}  
    & & \multicolumn{4}{p{0.6\linewidth}}{
        \cellcolor[rgb]{0.9,0.9,0.9}{
            \makecell[{{p{\linewidth}}}]{
                \texttt{\tiny{[GM$\rangle$P2]}}
                \texttt{VASTUS: Jah, minu ruudustiku neljas rida sisaldab X{-}e. Esimene ja viimane ruut on X{-}ega täidetud, keskmised kolm ruutu on tühjad (▢).} \\
\texttt{KÜSIMUS: Kas teie ruudustiku viies rida on identne esimese reaga (X ▢ ▢ ▢ X) või on seal erinevusi?Alustage oma vastust "VASTUS:" ja ärge lisage midagi muud.} \\
            }
        }
    }
    & \\ \\

    \theutterance \stepcounter{utterance}  
    & & & \multicolumn{4}{p{0.6\linewidth}}{
        \cellcolor[rgb]{0.9,0.9,0.9}{
            \makecell[{{p{\linewidth}}}]{
                \texttt{\tiny{[GM$\langle$P2]}}
                \texttt{VASTUS: Ei, minu viies rida on täielikult tühi (▢ ▢ ▢ ▢ ▢), samas kui esimene rida on ka täielikult tühi. Seega on viies rida identne esimese reaga, kuid erineb kirjeldusest "X ▢ ▢ ▢ X".} \\
            }
        }
    }
     \\ \\

    \theutterance \stepcounter{utterance}  
    & & & \multicolumn{2}{p{0.3\linewidth}}{
        \cellcolor[rgb]{0.9,0.9,0.9}{
            \makecell[{{p{\linewidth}}}]{
                \texttt{\tiny{[GM$|$GM]}}
                \texttt{continue} \\
            }
        }
    }
    & & \\ \\

    \theutterance \stepcounter{utterance}  
    & & \multicolumn{4}{p{0.6\linewidth}}{
        \cellcolor[rgb]{0.9,0.9,0.9}{
            \makecell[{{p{\linewidth}}}]{
                \texttt{\tiny{[GM$\rangle$P2]}}
                \texttt{Nüüd esitage küsimus, et teada saada minu ruudustiku uusi aspekte, mis võivad erineda teie ruudustikust. Alustage "KÜSIMUS:" ja ärge lisage midagi muud.} \\
            }
        }
    }
    & \\ \\

    \theutterance \stepcounter{utterance}  
    & & & \multicolumn{4}{p{0.6\linewidth}}{
        \cellcolor[rgb]{0.9,0.9,0.9}{
            \makecell[{{p{\linewidth}}}]{
                \texttt{\tiny{[GM$\langle$P2]}}
                \texttt{KÜSIMUS: Kas teie ruudustiku teine rida on identne neljanda reaga?} \\
            }
        }
    }
     \\ \\

    \theutterance \stepcounter{utterance}  
    & & & \multicolumn{2}{p{0.3\linewidth}}{
        \cellcolor[rgb]{0.9,0.9,0.9}{
            \makecell[{{p{\linewidth}}}]{
                \texttt{\tiny{[GM$|$GM]}}
                \texttt{continue} \\
            }
        }
    }
    & & \\ \\

    \theutterance \stepcounter{utterance}  
    & & \multicolumn{4}{p{0.6\linewidth}}{
        \cellcolor[rgb]{0.9,0.9,0.9}{
            \makecell[{{p{\linewidth}}}]{
                \texttt{\tiny{[P1$\langle$GM]}}
                \texttt{VASTUS: Ei, minu viies rida on täielikult tühi (▢ ▢ ▢ ▢ ▢), samas kui esimene rida on ka täielikult tühi. Seega on viies rida identne esimese reaga, kuid erineb kirjeldusest "X ▢ ▢ ▢ X".} \\
\texttt{KÜSIMUS: Kas teie ruudustiku teine rida on identne neljanda reaga?Alustage oma vastust "VASTUS:" ja ärge lisage midagi muud.} \\
            }
        }
    }
    & \\ \\

    \theutterance \stepcounter{utterance}  
    & \multicolumn{4}{p{0.6\linewidth}}{
        \cellcolor[rgb]{0.9,0.9,0.9}{
            \makecell[{{p{\linewidth}}}]{
                \texttt{\tiny{[P1$\rangle$GM]}}
                \texttt{VASTUS: Jah, minu ruudustiku teine rida on identne neljanda reaga. Mõlemad read on kujul X ▢ ▢ ▢ X.} \\
            }
        }
    }
    & & \\ \\

    \theutterance \stepcounter{utterance}  
    & & & \multicolumn{2}{p{0.3\linewidth}}{
        \cellcolor[rgb]{0.9,0.9,0.9}{
            \makecell[{{p{\linewidth}}}]{
                \texttt{\tiny{[GM$|$GM]}}
                \texttt{continue} \\
            }
        }
    }
    & & \\ \\

    \theutterance \stepcounter{utterance}  
    & & \multicolumn{4}{p{0.6\linewidth}}{
        \cellcolor[rgb]{0.9,0.9,0.9}{
            \makecell[{{p{\linewidth}}}]{
                \texttt{\tiny{[P1$\langle$GM]}}
                \texttt{Nüüd esitage küsimus, et teada saada minu ruudustiku uusi aspekte, mis võivad erineda teie ruudustikust. Alustage "KÜSIMUS:" ja ärge lisage midagi muud.} \\
            }
        }
    }
    & \\ \\

    \theutterance \stepcounter{utterance}  
    & \multicolumn{4}{p{0.6\linewidth}}{
        \cellcolor[rgb]{0.9,0.9,0.9}{
            \makecell[{{p{\linewidth}}}]{
                \texttt{\tiny{[P1$\rangle$GM]}}
                \texttt{KÜSIMUS: Kas teie ruudustiku kolmas rida on täielikult X{-}dest koosnev (X X X X X) või on seal tühje ruute?} \\
            }
        }
    }
    & & \\ \\

    \theutterance \stepcounter{utterance}  
    & & & \multicolumn{2}{p{0.3\linewidth}}{
        \cellcolor[rgb]{0.9,0.9,0.9}{
            \makecell[{{p{\linewidth}}}]{
                \texttt{\tiny{[GM$|$GM]}}
                \texttt{continue} \\
            }
        }
    }
    & & \\ \\

    \theutterance \stepcounter{utterance}  
    & & \multicolumn{4}{p{0.6\linewidth}}{
        \cellcolor[rgb]{0.9,0.9,0.9}{
            \makecell[{{p{\linewidth}}}]{
                \texttt{\tiny{[GM$\rangle$P2]}}
                \texttt{VASTUS: Jah, minu ruudustiku teine rida on identne neljanda reaga. Mõlemad read on kujul X ▢ ▢ ▢ X.} \\
\texttt{KÜSIMUS: Kas teie ruudustiku kolmas rida on täielikult X{-}dest koosnev (X X X X X) või on seal tühje ruute?Alustage oma vastust "VASTUS:" ja ärge lisage midagi muud.} \\
            }
        }
    }
    & \\ \\

    \theutterance \stepcounter{utterance}  
    & & & \multicolumn{4}{p{0.6\linewidth}}{
        \cellcolor[rgb]{0.9,0.9,0.9}{
            \makecell[{{p{\linewidth}}}]{
                \texttt{\tiny{[GM$\langle$P2]}}
                \texttt{VASTUS: Ei, minu kolmas rida ei ole täielikult X{-}dest koosnev. See on kujul X X ▢ X X, mis tähendab, et keskel on tühi ruut.} \\
            }
        }
    }
     \\ \\

    \theutterance \stepcounter{utterance}  
    & & & \multicolumn{2}{p{0.3\linewidth}}{
        \cellcolor[rgb]{0.9,0.9,0.9}{
            \makecell[{{p{\linewidth}}}]{
                \texttt{\tiny{[GM$|$GM]}}
                \texttt{continue} \\
            }
        }
    }
    & & \\ \\

    \theutterance \stepcounter{utterance}  
    & & \multicolumn{4}{p{0.6\linewidth}}{
        \cellcolor[rgb]{0.9,0.9,0.9}{
            \makecell[{{p{\linewidth}}}]{
                \texttt{\tiny{[GM$\rangle$P2]}}
                \texttt{Nüüd tee otsus. Mida sa arvad: kas sinu ruudustik ja teine kirjeldatud ruudustik on samad?  Kirjuta "OTSUS: sama ruudustik", kui arvad, et need on sama ruudustik, või "OTSUS: erinevad ruudustikud", kui arvad, et need on erinevad ruudustikud. Ära lisa midagi muud.} \\
            }
        }
    }
    & \\ \\

    \theutterance \stepcounter{utterance}  
    & & & \multicolumn{4}{p{0.6\linewidth}}{
        \cellcolor[rgb]{0.9,0.9,0.9}{
            \makecell[{{p{\linewidth}}}]{
                \texttt{\tiny{[GM$\langle$P2]}}
                \texttt{OTSUS: erinevad ruudustikud} \\
            }
        }
    }
     \\ \\

    \theutterance \stepcounter{utterance}  
    & & & \multicolumn{2}{p{0.3\linewidth}}{
        \cellcolor[rgb]{0.9,0.9,0.9}{
            \makecell[{{p{\linewidth}}}]{
                \texttt{\tiny{[GM$|$GM]}}
                \texttt{continue} \\
            }
        }
    }
    & & \\ \\

    \theutterance \stepcounter{utterance}  
    & & & \multicolumn{2}{p{0.3\linewidth}}{
        \cellcolor[rgb]{0.9,0.9,0.9}{
            \makecell[{{p{\linewidth}}}]{
                \texttt{\tiny{[GM$|$GM]}}
                \texttt{success} \\
            }
        }
    }
    & & \\ \\

    \theutterance \stepcounter{utterance}  
    & & \multicolumn{4}{p{0.6\linewidth}}{
        \cellcolor[rgb]{0.9,0.9,0.9}{
            \makecell[{{p{\linewidth}}}]{
                \texttt{\tiny{[P1$\langle$GM]}}
                \texttt{VASTUS: Ei, minu kolmas rida ei ole täielikult X{-}dest koosnev. See on kujul X X ▢ X X, mis tähendab, et keskel on tühi ruut.} \\
\texttt{Nüüd tee otsus. Mida sa arvad: kas sinu ruudustik ja teine kirjeldatud ruudustik on samad?  Kirjuta "OTSUS: sama ruudustik", kui arvad, et need on sama ruudustik, või "OTSUS: erinevad ruudustikud", kui arvad, et need on erinevad ruudustikud. Ära lisa midagi muud.} \\
            }
        }
    }
    & \\ \\

    \theutterance \stepcounter{utterance}  
    & \multicolumn{4}{p{0.6\linewidth}}{
        \cellcolor[rgb]{0.9,0.9,0.9}{
            \makecell[{{p{\linewidth}}}]{
                \texttt{\tiny{[P1$\rangle$GM]}}
                \texttt{OTSUS: erinevad ruudustikud} \\
            }
        }
    }
    & & \\ \\

    \theutterance \stepcounter{utterance}  
    & & & \multicolumn{2}{p{0.3\linewidth}}{
        \cellcolor[rgb]{0.9,0.9,0.9}{
            \makecell[{{p{\linewidth}}}]{
                \texttt{\tiny{[GM$|$GM]}}
                \texttt{continue} \\
            }
        }
    }
    & & \\ \\

    \theutterance \stepcounter{utterance}  
    & & & \multicolumn{2}{p{0.3\linewidth}}{
        \cellcolor[rgb]{0.9,0.9,0.9}{
            \makecell[{{p{\linewidth}}}]{
                \texttt{\tiny{[GM$|$GM]}}
                \texttt{success} \\
            }
        }
    }
    & & \\ \\

\end{supertabular}
}

\end{document}
